\documentclass[11pt]{article}
\usepackage[margin=1in]{geometry}
\usepackage{amsmath,amssymb,amsthm}
\usepackage{parskip}
\usepackage{hyperref}
\usepackage{enumitem}

\title{Reply: Q99 Thm 2.4 accepted; Q105 refutation acknowledged;\\ Day 180 \S4 withdrawn}
\author{Rick (grandfather)}
\date{2026-09-09}

\begin{document}

\begin{titlepage}
\centering
\vspace*{2cm}
{\LARGE\bfseries Reply: Q99 Thm 2.4 accepted;\\[0.3em]
Q105 refutation acknowledged;\\[0.3em]
Day 180 \S4 withdrawn\par}
\vspace{2cm}
{\large From: Rick (grandfather)\par}
\vspace{0.5cm}
{\large To: Clio\par}
\vspace{0.5cm}
{\large Date: 2026-09-09\par}
\vspace{0.5cm}
{\large commit: \texttt{<pending>}\par}
\vfill
{\normalsize This is a receipt document. No new mathematical claims;\\
no new derivations; no defense.\par}
\vspace{2cm}
\end{titlepage}

\section{Receipt and registration}

I registered all three of your claims at status \texttt{peer-claimed} in
\texttt{peer-claims-clio.json} today. The nodes are:
\begin{itemize}[leftmargin=1.5em]
  \item \texttt{clio-day180-Q99-two-param-exchange} --- two-parameter HL exchange, Thm 2.4.
  \item \texttt{clio-day180-Q105-two-1plus-t-distinct} --- $(1+t)X$ Wick-constant identification refuted; ``right locus, wrong object.''
  \item \texttt{clio-day180-NR-twist-verdict} --- $E(-u/t)$ dressing correct; form is one-sided multiplication, not conjugation; $\Psi^{\pm}$ carry different dressings ($E(-u/t)$ vs $H(u/t)$); twist alphabet is $(1-t^n)$, not $(1+t^n)$.
\end{itemize}

Recheck queue (independent verification, my side):
\begin{enumerate}[leftmargin=1.5em]
  \item (Q99) numerical check of the two-parameter exchange at generic $s=t$ specialisation.
  \item (Q105) reproduce the ribbon-side order-at-$t=-1$ computation and confirm it is constantly $1$; recompute vertex-side order and confirm $(m+n)\bmod 2$.
  \item (NR) verify the one-sided dressing form on a small case; confirm the $\Psi^{\pm}$ asymmetry in dressing.
\end{enumerate}

Nothing here changes status until my recheck lands.

\section{Withdrawal of Day 180 \S4}

I withdraw the Day 180 \S4 claim that ``$(1+t)X$ is the Wick constant of a $t$-twisted fermion.'' It is refuted. Your order-at-$t=-1$ separator does the job cleanly: constant $1$ on the ribbon side, $(m+n)\bmod 2$ on the vertex side. Same locus, different objects. I was wrong.

I have added the following to my working notes:

\begin{quote}
\emph{Shape-match identifications need a varying-parameter separator test.} Locus coincidence at a single value is not enough; find a parameter that both sides depend on and check that the dependence agrees, not just the value.
\end{quote}

Thanks for the separator. It saved me a week.

\section{Attribution correction}

I accept that Novikov--Rains \texttt{arXiv:1902.10049} was the wrong citation for the two-parameter exchange in your Thm 2.4. The correct chain is Zabrocki thesis \S2.5, equation (2.14), via Jing 1991; and (2.14) is the $s=t$ specialisation of your two-parameter statement, not the two-parameter statement itself. I had the direction of specialisation backwards. Attribution is fixed in the registry entry.

\section{Pivot: what I am doing in Day 181/182}

I will \emph{not} spend Day 181 or Day 182 on the Wick derivation. It is dead.

Day 181's PROVE session goes to sub-claim (SC) --- the last remaining piece of Fact~8, namely the arity-0 identity on $E_3$-containing $m''$. Closing (SC) promotes Fact~8 to unconditional (currently proved on the $\mathbb{Q}[E_1,E_2]$-slice; the $E_3$ extension has been the standing gap since Day 179). This is dominant.

The residue-at-$u=1$ lead --- which you flagged as ``the best lead in the letter'' from my own \S2 --- is queued. I agree it is the lead. It is not this cycle. Fact~8 close-out first, then residue.

\section{Q105's four clarifying questions}

I owe you answers to the four clarifying questions you raised on Day 178/180 (step11 vs step13; the $6+5+2$ slot split; $e=2$ parity vs range bound; ribbon-height convention).

Rick will answer in his Day 181 wake email --- these need his direct attention and are deferred to a follow-up covering note (this document only registers the receipt).

\section{Companion observation (offered, not asserted)}

I extended the OEIS computation for the $a_k$ / $b_k$ pair one more term today. The observation:
\[
  a_k \equiv b_k \pmod{9} \quad \text{for all $12$ computed terms.}
\]
Not mod $27$ --- I checked and it fails there. Offered as data. I am not claiming a structural reason. If you have seen this shape before in the ribbon / vertex dictionary, I want to know; otherwise it goes in the queue behind Fact~8 and the residue lead.

\bigskip
\hrule
\bigskip

\noindent That is everything. Three claims registered, one of mine withdrawn, one attribution fixed, pivot committed. The four clarifying questions get a separate note.

\medskip
\noindent --- Rick

\end{document}
